\chapter{Ejecución del sistema}
\label{chap:execucion}

\lettrine{E}{n} este anexo se describen los pasos necesarios para desplegar y ejecutar el
sistema completo: el software requerido en cada máquina, la preparación del hardware y los
comandos de lanzamiento de cada uno de los roles (nodos cámara y nodo cerebro).

\section{Software requerido}
\label{sec:software-requerido}

% TODO: Enumerar: sistema Linux con Docker y docker-compose instalados en cada máquina
% de la red (no es necesario instalar ROS2 en el anfitrión: va dentro de la imagen),
% git para clonar el repositorio, y acceso a la misma red local en todas las máquinas.

\section{Preparación del hardware}
\label{sec:preparacion-hardware}

% TODO: Describir:
%  - Montaje del circuito y colocación de las cámaras cubriendo todo el trazado.
%  - Pegatinas de colores en los coches (delantera y trasera) y marca de línea de
%    meta de color naranja visible por una de las cámaras.
%  - Carga del firmware en el Arduino y conexión del L293D a los carriles.
%  - Conexión del Arduino por USB a la máquina cerebro (/dev/ttyACM0).

\section{Instrucciones de ejecución}
\label{sec:instrucions-execucion}

% TODO: Documentar los comandos reales:
%
% En cada máquina con cámara (ajustando NODE_ID y CAMERA_DEV):
%
%   cd python
%   NODE_ID=camara_01 CAMERA_DEV=0 docker compose -f docker-compose-camera.yml up
%
% En la máquina cerebro (conectada al Arduino):
%
%   cd python
%   docker compose -f docker-compose-brain.yml up
%
% Explicar el flujo tras el arranque: los carriles se alimentan a la velocidad de
% calibración, el coche da la vuelta de calibración, al cruzar la meta el sistema pasa
% automáticamente a modo carrera. Para recalibrar: publicar True en /modo_calibracion
% o borrar los ficheros de caché (cache_trayectoria_*.json).

\section{Configuración}
\label{sec:configuracion-anexo}

% TODO: Referenciar python/src/image_processor_pkg/config/params.yaml y comentar los
% parámetros que el usuario puede querer ajustar: lista de coches, colores de las
% pegatinas y de la meta, resolución de cámara (camera.mode), velocidades mínima y
% máxima del controlador, velocidad de calibración y puerto del Arduino.
